% Modernised reading — fonds Alexandre Grothendieck, shelfmark 135
% (pages 1-58, the whole folder).
%
% This is an INTERPRETATION, not a transcription. It re-reads the pages in
% today's notation and vocabulary, states the hypotheses the manuscript leaves
% implicit, and drops the critical apparatus so that the mathematics reads
% straight through. Everything the brackets and underlines carried — a doubtful
% reading, a notation he used differently, a missing passage — moves into a
% footnote.
%
% It is held to being mathematically correct as it stands, not merely faithful
% to the page. Where the manuscript is loose or mistaken, the true statement is
% what appears here, and the footnote says what the page has. In any dispute of
% substance the transcription governs: whoever cites Grothendieck must cite
% batch-0{1,2,3}.fr.tex.
%
% Pass: Opus 5 (claude-opus-5), 2026-08-16 — first pass, unchecked against the
% pages by a human. The folder was read whole in one pass, all three batches
% (pages 1-58) together, and this is a SINGLE document for the whole folder
% rather than one file per batch. Two reasons, both of them the folder's own.
% Hoàng Xuân Sính's typescript starts on page 39 and ends on page 44: the batch
% boundary falls inside one of its sentences, and three files would have had to
% cut it there. And the folder tells one argument three times over, in three
% hands — Grothendieck's notes, the typescript, the English manuscript — so
% that almost every section of this reading draws on two batches at once.
%
% One half of the folder's source is in English (pages 46-58, a thirteen-page
% manuscript). This edition is French throughout, so that half is restated
% across a language boundary rather than lightly edited; the résumé says so.
\documentclass[11pt,a4paper]{article}
% Le préambule commun à toutes les transcriptions du fonds Grothendieck.
%
% Il tient en une idée : **l'incertitude se compose, elle ne se résout pas.**
% Une transcription qui lisse un mot douteux a détruit la seule chose pour
% laquelle on la faisait — un lecteur doit pouvoir savoir, à chaque endroit,
% ce qui était sur le feuillet et ce qui a été ajouté. Les six macros qui
% suivent sont donc l'appareil critique, et non de la décoration.
%
% Les mêmes commandes sont reconnues par `scripts/render.mjs`, qui produit la
% vue de lecture du site. Une macro ajoutée ici doit l'être aussi là-bas,
% faute de quoi le rendu échouera bruyamment — ce qui est voulu : un rendu
% silencieusement faux est pire qu'un rendu absent.

\ProvidesPackage{grothendieck}[2026/08/08 Transcriptions du fonds Grothendieck]

\usepackage{amsmath,amssymb,amsthm}
\usepackage{mathrsfs}
\usepackage{stmaryrd}
% Les diagrammes commutatifs sont partout dans ce fonds : c'est la langue
% même de la géométrie algébrique de Grothendieck, et une transcription qui
% les rendrait en prose ne transcrirait rien.
\usepackage{tikz-cd}
% Français par défaut — c'est la langue des feuillets — mais l'anglais est
% chargé aussi : la traduction et le résumé sont des documents anglais, et la
% typographie française (espaces avant les deux-points) y serait une faute.
% Ces documents basculent par \selectlanguage{english} en tête de corps.
\usepackage[english,french]{babel}
\usepackage{fontspec}
\usepackage{geometry}
\geometry{a4paper,margin=25mm,marginparwidth=28mm}
\usepackage{marginnote}
\usepackage{xcolor}
% ulem, not soul. `soul`'s \st and \ul are built on a character-by-character
% scan that XeLaTeX's font machinery breaks: the document fails with an
% undefined control sequence rather than degrading. ulem does the same two jobs
% and is Unicode-engine safe. `normalem` stops it hijacking \emph, which the
% transcriptions use for Grothendieck's own emphasis.
\usepackage[normalem]{ulem}
\usepackage{hyperref}
\hypersetup{colorlinks=true,linkcolor=black,urlcolor=black,pdfborder={0 0 0}}

% --- Métadonnées du lot -----------------------------------------------------
% Reprises telles quelles par le script de rendu, qui les affiche en tête de
% la vue de lecture. Elles fixent ce que le fichier prétend couvrir : sans
% elles, une transcription flotte sans point d'attache dans un fonds de seize
% mille feuillets.

\newcommand{\folder}[1]{\gdef\grothFolder{#1}}
\newcommand{\batch}[1]{\gdef\grothBatch{#1}}
\newcommand{\leaves}[2]{\gdef\grothFirst{#1}\gdef\grothLast{#2}}
\newcommand{\foldertitle}[1]{\gdef\grothFoldertitle{#1}}
\gdef\grothFolder{?}\gdef\grothBatch{?}\gdef\grothFirst{?}\gdef\grothLast{?}\gdef\grothFoldertitle{}

% --- Le résumé --------------------------------------------------------------
%
% Il ouvre la lecture modernisée, sous le titre « Résumé », et il est composé
% en petit. Ce n'est pas un ornement : c'est le seul passage du document qui
% ne soit pas des mathématiques, et un lecteur doit pouvoir le reconnaître
% avant de l'avoir lu — puis le sauter s'il connaît déjà le sujet, ou s'y
% arrêter s'il ne le connaît pas. Le filet qui le suit marque la reprise du
% corps.

\newenvironment{resume}
  {\small\setlength{\parskip}{0.45em}}
  {\par\medskip\hrule\medskip}

% --- Mots-clés ---------------------------------------------------------------
%
% En anglais, en clôture du résumé de la lecture modernisée : le vocabulaire
% moderne sous lequel chercher ce que les feuillets construisent. Le manifeste
% du site (`npm run manifest`) les extrait pour étiqueter la cote entière —
% cette ligne est la source unique des tags, il n'y a pas de fichier à côté.
\newcommand{\keywords}[1]{\par\smallskip\noindent
  {\footnotesize\textsc{Keywords} --- \textit{#1}}\par}

% --- L'appareil critique ----------------------------------------------------

% Le numéro de feuillet, en marge. C'est l'ancre qui permet de régler un
% désaccord sur une formule en regardant *un* feuillet plutôt qu'une chemise
% de six cents. Le site s'en sert aussi pour tourner les pages du fac-similé
% quand on fait défiler la transcription.
\newcommand{\leaf}[1]{%
  \marginnote{\footnotesize\color{gray!70}\textsf{#1}}%
  \label{leaf:#1}\ignorespaces}

% Illisible : on ne devine pas. Un mot inventé qui se lit comme les autres est
% le pire résultat possible.
\newcommand{\ill}{\textcolor{red!60!black}{[\,\ldots\,]}}

% Lecture douteuse : le mot est donné, et signalé comme incertain.
\newcommand{\uncertain}[1]{\uline{#1}}

% Ajout de l'éditeur — une lettre manquante, un mot sous-entendu.
\newcommand{\add}[1]{\textcolor{blue!50!black}{[#1]}}

% Biffé par Grothendieck lui-même. Ce qu'il a rayé fait partie du document :
% une rature dit souvent où la pensée a changé de direction.
\newcommand{\struck}[1]{\sout{#1}}

% Note du transcripteur, sur l'état matériel du feuillet.
\newcommand{\note}[1]{\par\smallskip{\small\color{gray!60!black}\textit{[#1]}}\par\smallskip}

% Note marginale de Grothendieck, distincte du corps du texte.
\newcommand{\marginal}[1]{\par\smallskip\noindent\hspace{1em}%
  {\small\color{gray!60!black}#1}\par\smallskip}

% --- Titre ------------------------------------------------------------------

\AtBeginDocument{%
  \begin{center}
    {\large\bfseries Fonds Alexandre Grothendieck --- cote n\textsuperscript{o}~\grothFolder}\\[2pt]
    {\itshape\grothFoldertitle}\\[4pt]
    {\small feuillets \grothFirst--\grothLast\ (lot \grothBatch)}\\[2pt]
    {\footnotesize\color{gray!60!black}Transcription établie d'après le fac-similé numérisé
     par l'Université de Montpellier.\\
     Les lectures incertaines sont soulignées, les illisibles notées [\,\ldots\,],
     les ajouts entre crochets.}
  \end{center}
  \vspace{1em}}

\endinput


\folder{135}
\pages{1}{58}
\dating{[vers 1974]}
\watermark{Édition de démonstration\\interprétation personnelle de l'œuvre}
\foldertitle{Catégories [et Gr-catégories] : notes manuscrites (s.d.), copies de tapuscrit annoté (s.d.), copies de manuscrit annoté (s.d.) — lecture modernisée du dossier entier}

\begin{document}

\section*{Résumé}

\begin{resume}
Ces pages demandent ce que devient la notion de groupe lorsqu'on cesse
d'exiger que la multiplication soit associative, et qu'on demande seulement
qu'elle le soit à un isomorphisme près — un isomorphisme qu'il faut alors
choisir, et qui devient à son tour l'objet d'étude.

Dans un groupe, $(xy)z$ et $x(yz)$ sont le même élément et il n'y a rien à
dire. Supposons maintenant que les « éléments » soient les objets d'une
catégorie et que l'égalité ait été remplacée par l'isomorphisme : il faut
alors se donner, pour chaque triplet d'objets, un isomorphisme précis entre
les deux parenthésages, et vérifier que les différentes manières de
réassocier quatre facteurs s'accordent. Cette donnée n'est pas canonique.
Deux choix diffèrent, et c'est cette différence — non le groupe qu'on
retrouve en oubliant tout — qui est le sujet du dossier.

L'idée qui arrive est une idée d'obstruction, et elle a un précédent bien
connu. Une extension d'un groupe $G$ par un groupe commutatif $M$ est
classifiée par une classe dans $H^{2}(G,M)$ : la classe mesure ce qui empêche
l'extension d'être triviale, et toutes les classes se produisent. Ici la même
chose arrive un cran plus haut. On choisit un représentant $L_{a}$ pour chaque
classe d'objets et un isomorphisme $L_{a} \otimes L_{b} \simeq L_{ab}$ ;
l'associativité devient alors une fonction $f(a,b,c)$ de trois variables, à
valeurs dans le groupe des automorphismes de l'objet unité ; l'axiome de
cohérence dit exactement que $f$ est un $3$-cocycle ; changer les choix change
$f$ d'un cobord. Il reste une seule classe
$k(C) \in H^{3}(\pi_{0}(C), \pi_{1}(C))$, et elle détermine tout.

Une seconde couche apparaît si la multiplication est aussi commutative à
isomorphisme près. L'invariant supplémentaire est minuscule : c'est ce que
fait la symétrie appliquée à $X \otimes X$, un élément d'ordre $2$. En algèbre
graduée c'est la règle des signes de Koszul — une droite de degré impair,
échangée avec elle-même, donne $-1$ — et ce seul signe est tout l'invariant.
Le dossier en contient l'exemple, à la page 32, et c'est le plus petit
possible.

Où cela mène-t-il ? Une catégorie de Picard est aujourd'hui la même chose
qu'un spectre connectif tronqué en degré $1$ ; la classe $k(C)$ est son
premier invariant de Postnikov, et le signe ci-dessus est la multiplication
par l'élément $\eta$ de la sphère. Et la K-théorie est au bout : l'« enveloppe
de Picard » de la catégorie des modules projectifs de type fini sur un anneau
$R$ a pour invariants $K^{0}(R)$ et $K^{1}(R)$, ce que le dossier énonce, et
il conjecture la version en dimension $n$, avec les $K_{i}$ pour invariants —
ce qui est le théorème de complétion en groupe démontré, dans ces mêmes
années, par Segal et par May.

Le dossier est fait de trois mains. Les pages 2 à 37 sont des notes
manuscrites ; les pages 39 à 44, la copie annotée d'un tapuscrit de Hoàng Xuân
Sính présenté par Henri Cartan ; les pages 46 à 58, un manuscrit de treize
pages rédigé en anglais, qui donne le plan d'un travail en trois chapitres —
la thèse que le tapuscrit annonce. Ce n'est pas trois sujets mais un seul
argument raconté deux fois, et c'est ce qui rend la lecture d'ensemble
possible : on peut mettre côte à côte l'énoncé d'un théorème et la note qui
l'esquisse. Entre les deux, du papier de récupération — une chanson, un poème,
et la feuille de calendrier hanoïen de 1974 dont les archivistes ont
vraisemblablement tiré la datation du dossier.

La dernière page mathématique est une prophétie, et elle n'a besoin d'aucune
connaissance préalable : algèbre homologique, algèbre homotopique et algèbre
catégorique doivent fusionner en une seule vision, à quoi conviendrait le nom
d'algèbre homologique non commutative, « ce qui porte actuellement ce nom
n'[étant] qu'une amorce très partielle ». Les noms modernes à chercher sont
\emph{2-groupe}, \emph{théorème de Sính}, \emph{groupoïde de Picard},
\emph{2-gerbe} et \emph{enveloppe de Picard}.

\keywords{2-group, Sinh's theorem, Picard groupoid, 2-gerbe, Picard envelope,
algebraic K-theory}
\end{resume}

\pagerange{2}{58}
\section*{Le dossier, ses trois mains, et l'ordre suivi ici}

Le dossier ne se lit pas dans l'ordre où il est relié. Les notes des pages 2 à
37 sont d'un travail en cours et sautent d'un sujet à l'autre ; le tapuscrit
des pages 39 à 44 et le manuscrit anglais des pages 46 à 58 exposent, eux, une
théorie déjà faite. La présente lecture regroupe par argument : elle prend les
définitions là où elles sont complètes — dans le manuscrit anglais — puis le
théorème de classification, puis ses spécialisations, puis les constructions
qui en dépendent, et rapproche donc régulièrement des pages éloignées les unes
des autres. Le repère de page en marge de chaque section dit d'où elle est
tirée.\footnote{Les pages 1, 6, 24, 30, 34, 38 et 45 sont blanches. Les pages 3
et 5 portent deux copies d'une chanson anglaise, la page 33 un poème français
biffé, la page 36 une feuille de calendrier de 1974 émise à Hanoï — du papier
de récupération, dont la dernière feuille est vraisemblablement à l'origine de
la datation « [vers 1974] » de l'inventaire.}

Trois conventions sont fixées ici pour tout le dossier, parce qu'il en emploie
plusieurs. Les invariants sont notés $\pi_{0}$ et $\pi_{1}$, comme dans les
notes et le tapuscrit, et non $\Pi_{0}$, $\Pi_{1}$ comme dans le manuscrit
anglais. L'homomorphisme qui mesure le défaut de stricte commutativité est
noté $s$, comme dans le tapuscrit ; les notes l'écrivent $\sigma$, et parfois
$\Sigma_{\varphi}$ lorsqu'il provient d'une forme bilinéaire $\varphi$. Enfin
la contrainte d'associativité est prise dans le sens usuel aujourd'hui,
$a_{X,Y,Z} : (X \otimes Y) \otimes Z \to X \otimes (Y \otimes Z)$, alors que
le manuscrit anglais l'écrit dans l'autre sens.\footnote{Page 46 :
$a_{X,Y,Z} : X \otimes (Y \otimes Z) \xrightarrow{\ \sim\ } (X \otimes Y)
\otimes Z$. Le choix est sans conséquence — les contraintes sont inversibles —
mais il retourne tous les pentagones et tous les hexagones du manuscrit, et il
vaut mieux le dire une fois que le retrouver à chaque diagramme.}

Le vocabulaire du dossier est conservé. Une \emph{Gr-catégorie} est ce qu'on
appelle aujourd'hui un \emph{2-groupe}, ou groupe catégorique ; une
\emph{catégorie de Picard} est un \emph{groupoïde de Picard}, c'est-à-dire un
2-groupe symétrique. On garde les premiers noms, qui sont ceux du dossier et
de son titre, en rappelant les seconds là où ils servent à chercher.

\pagerange{46}{50}
\section*{Une Gr-catégorie : un groupe dont l'égalité est devenue un
isomorphisme}

Le manuscrit anglais des pages 46 à 50 est le seul endroit du dossier où les
définitions soient données au complet ; on les prend là.\footnote{Elles y sont
attribuées à la thèse de Neantro Saavedra Rivano, dont la terminologie est
suivie. Le tapuscrit de la page 39 renvoie à la même source.}

Une $\otimes$-catégorie est une catégorie $C$ munie d'un bifoncteur covariant
$\otimes : C \times C \to C$. Une \emph{contrainte d'associativité} est un
isomorphisme de trifoncteurs
\[
a_{X,Y,Z} : (X \otimes Y) \otimes Z \xrightarrow{\ \sim\ }
X \otimes (Y \otimes Z)
\]
satisfaisant l'\emph{axiome du pentagone} : pour tous $X, Y, Z, T$, le
diagramme

\begin{tikzcd}[column sep=tiny, row sep=small, nodes={font=\scriptsize}]
& ((X \otimes Y) \otimes Z) \otimes T
\arrow[dl, "a_{X \otimes Y{,}Z{,}T}"'] \arrow[dr, "a_{X{,}Y{,}Z} \otimes \mathrm{id}_{T}"] & \\
(X \otimes Y) \otimes (Z \otimes T) \arrow[d, "a_{X{,}Y{,}Z \otimes T}"'] & &
(X \otimes (Y \otimes Z)) \otimes T \arrow[d, "a_{X{,}Y \otimes Z{,}T}"] \\
X \otimes (Y \otimes (Z \otimes T)) & &
X \otimes ((Y \otimes Z) \otimes T) \arrow[ll, "\mathrm{id}_{X} \otimes a_{Y{,}Z{,}T}"]
\end{tikzcd}

commute. Une \emph{contrainte de commutativité} est un isomorphisme de
bifoncteurs $c_{X,Y} : X \otimes Y \xrightarrow{\sim} Y \otimes X$ tel que
$c_{Y,X} \circ c_{X,Y} = \mathrm{id}_{X \otimes Y}$ ; elle est dite
\emph{stricte} si $c_{X,X} = \mathrm{id}_{X \otimes X}$ pour tout $X$. Une
\emph{contrainte d'unité} est un triplet $(1, g, d)$ où $1$ est un objet et
$g_{X} : X \xrightarrow{\sim} 1 \otimes X$, $d_{X} : X \xrightarrow{\sim} X
\otimes 1$ sont des isomorphismes naturels avec $g_{1} = d_{1}$. La
compatibilité de $a$ et $c$ est l'\emph{axiome de l'hexagone}, celle de $a$ et
de l'unité l'axiome du triangle ; une $\otimes$-catégorie munie de $a$ et de
l'unité est dite AU, munie des trois ACU.

Un \emph{$\otimes$-foncteur} $C \to C'$ est un couple $(F, \check{F})$ formé
d'un foncteur $F$ et d'un isomorphisme de bifoncteurs $\check{F}_{X,Y} : FX
\otimes FY \to F(X \otimes Y)$ ; il est dit associatif, commutatif, unifère
selon qu'il respecte $a$, $c$, l'unité — l'isomorphisme $\hat{F} : 1' \simeq
F1$ d'un foncteur unifère étant unique lorsqu'il existe. Un
\emph{$\otimes$-morphisme} $(F,\check{F}) \to (G,\check{G})$ est un morphisme
de foncteurs $\lambda$ compatible aux $\check{F}$, $\check{G}$.

\textbf{Définition.} Une \emph{Gr-catégorie} est une $\otimes$-catégorie AU
dont la catégorie sous-jacente est un groupoïde et dont tout objet est
inversible.\footnote{« Inversible » : il existe $X'$, $X''$ avec $X' \otimes X
\simeq X \otimes X'' \simeq 1$. Dans un groupoïde monoïdal les deux inverses
coïncident à isomorphisme près, de sorte que la définition du manuscrit
(page 48) et la définition usuelle, qui demande un seul inverse bilatère, se
valent.} Une \emph{catégorie de Picard} est une Gr-catégorie munie d'une
contrainte de commutativité compatible avec sa contrainte d'associativité ;
elle est \emph{stricte} si cette contrainte l'est. Une Gr-catégorie est donc,
mot pour mot, un groupe où l'égalité a été remplacée par l'isomorphisme, et
une catégorie de Picard un groupe commutatif au même sens.

À une Gr-catégorie $C$ on attache trois invariants.

\begin{enumerate}
\item Le groupe $\pi_{0}(C)$ des classes d'isomorphisme d'objets, la loi étant
induite par $\otimes$.
\item Le groupe $\pi_{1}(C) = \mathrm{Aut}(1_{C})$, qui est
\emph{commutatif}.\footnote{C'est l'argument d'Eckmann et Hilton : sur
$\mathrm{Aut}(1)$ la composition et le produit $\otimes$ sont deux lois
unifères qui se distribuent l'une l'autre, donc elles coïncident et sont
commutatives. Le tapuscrit (page 40) et le manuscrit anglais (page 49)
énoncent le fait sans le démontrer.}
\item Une action de $\pi_{0}(C)$ sur $\pi_{1}(C)$, qui en fait un
$\pi_{0}(C)$-module à gauche. Pour $X$ un objet, les deux isomorphismes
$\gamma_{X} : u \mapsto u \otimes \mathrm{id}_{X}$ et $\delta_{X} : u \mapsto
\mathrm{id}_{X} \otimes u$ de $\mathrm{Aut}(1)$ sur $\mathrm{Aut}(X)$ ne
coïncident pas en général, et l'action de la classe $s$ de $X$ est
$s \cdot u = \delta_{X}^{-1}\gamma_{X}(u)$.
\end{enumerate}

Un \emph{épinglage} de type $(M,N)$ — $M$ un groupe, $N$ un $M$-module à
gauche — est un couple d'isomorphismes $\varepsilon_{0} : M \xrightarrow{\sim}
\pi_{0}(C)$, $\varepsilon_{1} : N \xrightarrow{\sim} \pi_{1}(C)$ compatibles
aux actions.\footnote{Le manuscrit anglais dit \emph{preepinglage} et
\emph{preepingled}, francisme resté dans un texte anglais ; les notes de
Grothendieck disent « épinglée par $(G,M)$ », parfois « d'invariants $G$,
$M$ ». On dit ici épinglage.} Une flèche entre Gr-catégories épinglées de même
type est un $\otimes$-foncteur associatif compatible aux épinglages ; une
telle flèche est automatiquement une $\otimes$-équivalence, de sorte que
classer ces objets à équivalence près, c'est compter les composantes connexes
de la catégorie qu'ils forment.

\pagerange{40}{41}
\section*{Le théorème de classification}

Soit $C$ une Gr-catégorie d'invariants $\pi_{0}$, $\pi_{1}$. Choisissons pour
chaque $a \in \pi_{0}$ un représentant $L_{a} \in \mathrm{Ob}\,C$, et pour
chaque couple un isomorphisme $\phi_{a,b} : L_{a} \otimes L_{b} \simeq
L_{ab}$. Pour trois éléments $a$, $b$, $c$, la contrainte d'associativité
\[
(L_{a} \otimes L_{b}) \otimes L_{c} \xrightarrow{\ \sim\ }
L_{a} \otimes (L_{b} \otimes L_{c}),
\]
transportée par $\phi_{a,b}$, $\phi_{ab,c}$, $\phi_{b,c}$ et $\phi_{a,bc}$,
devient un automorphisme de $L_{abc}$, c'est-à-dire la tensorisation par un
élément bien déterminé $f(a,b,c) \in \pi_{1}$. Les $\phi$ étant choisis, la
donnée de la contrainte d'associativité équivaut donc à celle d'une
application
\[
f : \pi_{0} \times \pi_{0} \times \pi_{0} \longrightarrow \pi_{1},
\]
et l'axiome du pentagone dit exactement que $f$ est un \emph{3-cocycle} de
$\pi_{0}$ à valeurs dans $\pi_{1}$. L'indétermination sur le système
$(\phi_{a,b})$ est celle d'une 2-cochaîne arbitraire $g$, et remplacer $\phi$
par $\phi g$ remplace $f$ par $f + \delta g$. Il reste une classe bien
définie
\[
k(C) \in H^{3}\bigl(\pi_{0}(C),\, \pi_{1}(C)\bigr).
\]

\textbf{Théorème} (Sính). \emph{L'application $C \mapsto k(C)$ est une
bijection de l'ensemble des classes d'équivalence de Gr-catégories épinglées
de type $(M,N)$ sur $H^{3}(M,N)$. Toute classe se produit, et une
Gr-catégorie est donc déterminée, à équivalence près, par le triplet formé du
groupe $M$, du $M$-module $N$ et d'une classe de cohomologie arbitraire dans
$H^{3}(M,N)$.}\footnote{Énoncé page 41 pour le tapuscrit, page 50 pour le
manuscrit anglais ; le second dit « canonical bijection ». L'épinglage n'est
pas une commodité : sans lui il faudrait encore quotienter par les
automorphismes de $M$ et de $N$, et l'ensemble des classes ne serait plus un
groupe. C'est le théorème auquel le nom de Hoàng Xuân Sính est aujourd'hui
attaché.}

La loi de groupe de $H^{3}$ a de plus une interprétation géométrique, à la
Baer : elle correspond à une opération sur les Gr-catégories
elles-mêmes.\footnote{Page 41. La construction n'est pas explicitée là ; les
notes des pages 11 et 17 s'en servent sous le nom de « multiplication de
Baer », et c'est elle qui fait des 2-catégories de la page 17 des objets
munis d'une loi.}

\pagerange{39}{41}
\section*{Ce que la classe $k(C)$ redonne dans les exemples}

Le tapuscrit accompagne le théorème de quatre exemples, et c'est là qu'on voit
ce qu'il vaut : deux d'entre eux redonnent des invariants classiques, les deux
autres en produisent de nouveaux.

\textbf{Les lacets d'un espace pointé.} Soit $(X,x)$ un espace topologique
pointé, $C$ la catégorie dont les objets sont les lacets de $X$ en $x$ et les
morphismes les classes d'homotopie d'homotopies entre lacets, $\otimes$ étant
la composition des lacets — associative seulement à homotopie près, ce qui est
tout le point. Alors $\pi_{0}(C) = \pi_{1}(X,x)$ et $\pi_{1}(C) =
\pi_{2}(X,x)$, et $k(C)$ est le premier invariant de Postnikov de $X$, dans
$H^{3}(\pi_{1}(X), \pi_{2}(X))$.\footnote{Le tapuscrit dit « on doit
retrouver », le mot « doit » étant une correction à l'encre sur la page 41 :
l'identification est annoncée, non démontrée. C'est l'exemple qui montre que
le théorème ne produit pas un invariant nouveau mais une nouvelle
interprétation d'un invariant connu — d'où la remarque, sur la même page,
d'une interprétation « géométrique » nouvelle. La variante donnée juste avant,
où $C$ est attachée à un espace de Hopf $E$ homotopiquement associatif,
contient celle-ci en prenant pour $E$ l'espace des lacets.}

\textbf{Les torseurs sous un faisceau abélien.} Si $F$ est un faisceau abélien
sur un topos $X$, la catégorie des torseurs sous $F$ munie de la composition
de Baer est une catégorie de Picard stricte, d'invariants $\pi_{0} =
H^{1}(X,F)$ et $\pi_{1} = H^{0}(X,F)$, et $k(C) = 0$. Les Modules inversibles
sur un topos annelé $(X, O_{X})$ en sont le cas $F = O_{X}^{*}$.

\textbf{Les auto-équivalences d'une catégorie.} Si $A$ est une catégorie, la
sous-catégorie pleine de $\mathrm{Hom}(A,A)$ formée des équivalences de $A$
avec elle-même, munie de la composition, est une Gr-catégorie ; $\pi_{0}$ est
le groupe des classes d'isomorphisme d'auto-équivalences, $\pi_{1}$ le groupe
des automorphismes du foncteur identique. Ici $k(C)$ « semble nouveau », dit
le tapuscrit, y compris lorsque $A$ est la catégorie des bi-Modules
inversibles sous un Anneau $O_{X}$. On peut prendre pour $A$, plus
généralement, un objet d'une 2-catégorie quelconque.

\textbf{Les bitorseurs, et l'invariant de Mac Lane.} Soit $F$ un faisceau en
groupes, non nécessairement commutatifs. Prenant pour $A$ le champ des
$F$-torseurs à droite, la Gr-catégorie de ses auto-équivalences s'identifie à
celle des \emph{bitorseurs} sous $F$ — les faisceaux munis de deux actions de
$F$ qui commutent et font chacune de $P$ un torseur — la loi étant la
composition de Baer. Ses invariants sont
\[
\pi_{0} = \mathrm{Autext}(F), \qquad \pi_{1} = Z(F),
\]
le groupe des automorphismes extérieurs et le centre.\footnote{La page 10
écrit $\pi_{0}(P(F)) \in \mathrm{Autext}\,F$ ; le signe voulu est l'égalité, ce
que la transcription signale. $\mathrm{Autext} = \mathrm{Out}$.} Et $k(C)$
n'est autre que l'invariant classique de Mac Lane, l'obstruction à réaliser
par une extension de groupes le « noyau » que $F$ définit.

\pagerange{41}{42}
\section*{Le cas commutatif : les catégories de Picard}

Soit $C$ une catégorie de Picard. Comme Gr-catégorie elle a ses invariants
$\pi_{0}$, $\pi_{1}$, mais $\pi_{0}$ est ici commutatif et agit trivialement
sur $\pi_{1}$ ; c'est une condition nécessaire pour qu'une Gr-catégorie porte
une structure de Picard. À tout objet $X$ on associe la symétrie de $X \otimes
X$, qui est un automorphisme de $X \otimes X$, donc la tensorisation par un
élément $s(X)$ de $\pi_{1}(C)$, et l'on obtient un homomorphisme canonique
\[
s = s(C) : \pi_{0} \longrightarrow {}_{2}\pi_{1},
\]
où ${}_{2}\pi_{1}$ est le sous-groupe des éléments annulés par $2$.

\textbf{Théorème} (Sính). \emph{Une catégorie de Picard est déterminée, à
équivalence non canonique de catégories de Picard près, par la donnée des
groupes commutatifs $\pi_{0}$ et $\pi_{1}$ et de l'homomorphisme $s : \pi_{0}
\to {}_{2}\pi_{1}$. Elle est stricte si et seulement si $s = 0$, et deux
catégories de Picard strictes de mêmes invariants $\pi_{0}$, $\pi_{1}$ sont
équivalentes.}\footnote{Page 42 ; « non canonique » est une addition à l'encre
sur la page, et elle compte : il n'y a pas d'équivalence privilégiée entre
deux catégories de Picard de mêmes invariants, seulement une équivalence. Le
cas strict est le résultat de Deligne, SGA 4 XVIII, auquel le tapuscrit
renvoie. Les notes disent la même chose en d'autres termes page 11 : la
classification des catégories de Picard épinglées par $(M,N)$ est, comme
groupe, $\mathrm{Hom}(M, {}_{2}N)$.}

Une conséquence mérite d'être isolée, parce qu'elle se perd facilement : une
équivalence de Gr-catégories entre deux catégories de Picard n'est pas
nécessairement compatible avec les contraintes de commutativité, donc n'est
pas nécessairement une équivalence de catégories de Picard. La classification
« au sens Picard » est strictement plus fine que celle des Gr-catégories.

\pagerange{11}{12}
\subsection*{La contrainte de commutativité comme torseur}

Les notes des pages 11 et 12 disent d'où vient l'invariant $s$, et la réponse
est qu'il est le seul possible. Fixons la contrainte d'associativité et
demandons quelles contraintes de commutativité lui sont compatibles. Si $c$ et
$c'$ en sont deux, leur différence est une application $\varphi : \pi_{0}
\times \pi_{0} \to \pi_{1}$ ; l'axiome de l'hexagone force $\varphi$ à être
bilinéaire, et la relation $c_{Y,X}c_{X,Y} = \mathrm{id}$ force $\varphi(x,y)
+ \varphi(y,x) = 0$. Réciproquement toute telle $\varphi$ tord $c$ en une
contrainte compatible. Donc :

\begin{itemize}
\item les structures de Picard sur une Gr-catégorie donnée forment un torseur
sous le groupe $\mathrm{Bil}^{\mathrm{as}}(\pi_{0}, \pi_{0}; \pi_{1})$ des
applications bilinéaires antisymétriques ;
\item les structures de Picard \emph{strictes} en forment un sous-torseur sous
le sous-groupe $\mathrm{Bil}^{\mathrm{alt}}(\pi_{0}, \pi_{0}; \pi_{1})$ des
applications bilinéaires \emph{alternées}, celles pour lesquelles
$\varphi(x,x) = 0$.\footnote{La page 12 écrit deux mots que la transcription
lit « Bilant » — ou « Bilalt », le premier étant surchargé d'une correction —
et « Bilalt ». Les deux énoncés ne sont distincts que si les deux mots le
sont, et la lecture retenue ici est la seule qui les rende tous deux vrais :
antisymétrique pour les structures de Picard, alternée pour les strictes. Sur
un groupe où $2$ n'est pas inversible, alterné est strictement plus fort
qu'antisymétrique.}
\end{itemize}

L'application $\varphi \mapsto (x \mapsto \varphi(x,x))$ envoie
$\mathrm{Bil}^{\mathrm{as}}$ dans $\mathrm{Hom}(\pi_{0}, {}_{2}\pi_{1})$, de
noyau $\mathrm{Bil}^{\mathrm{alt}}$, et c'est elle qui fait varier $s$. Tout
tient donc à savoir si elle est surjective. Elle l'est :

\textbf{Lemme} (page 4). \emph{Soient $M$, $N$ deux groupes commutatifs et
$\Sigma : M \to {}_{2}N$ un homomorphisme. Il existe une application
bilinéaire antisymétrique $\varphi : M \times M \to N$ telle que
$\varphi(x,x) = \Sigma(x)$ pour tout $x$.}

En effet $\Sigma$ est à valeurs dans ${}_{2}N$ et annule donc $2M$ : il se
factorise en une application $\mathbf{F}_{2}$-linéaire $M/2M \to {}_{2}N$.
Choisissons une base $(e_{i})_{i \in I}$ du $\mathbf{F}_{2}$-espace vectoriel
$M/2M$, posons $n_{i} = \Sigma(e_{i})$, et
\[
\varphi(e_{i}, e_{j}) =
\begin{cases}
0 & \text{si } i \neq j,\\
n_{i} & \text{si } i = j,
\end{cases}
\]
prolongée par bilinéarité. Alors $\varphi(x,y) + \varphi(y,x) = 2\varphi(x,y)
= 0$ puisque les valeurs sont dans ${}_{2}N$, et $\varphi(x,x) = \sum_{i}
x_{i}^{2} n_{i} = \sum_{i} x_{i} n_{i} = \Sigma(x)$ puisque $x_{i}^{2} =
x_{i}$ dans $\mathbf{F}_{2}$.\footnote{La page écrit « remplaçant $M$, $N$
resp. par $M_{2}$, ${}_{2}N$, on peut supposer $M$, $N$ des vectoriels sur
$\mathbf{F}_{2}$ » ; la transcription note que la lettre qui précède $M_{2}$
est indistincte. C'est $M/2M$ qu'il faut, et c'est ce qui est écrit ici : le
quotient, non le sous-groupe des éléments de $2$-torsion. La page conclut sur
un « OK » de sa main. Le choix d'une base est un choix : sur un topos il n'y
en a pas toujours, et c'est exactement ce que les pages 15 et 16 reprennent.}

\textbf{Corollaire} (page 4). \emph{Toute catégorie de Picard devient stricte
en changeant sa seule contrainte de commutativité. En particulier, comme
Gr-catégorie elle est splittée : $k(C) = 0$ pour toute catégorie de
Picard.}\footnote{Le second membre est, sur la page, un ajout interlinéaire
serré : « en tant que Gr-catégorie, elle est de type triviale ». Il se déduit
du premier : après torsion on a une catégorie de Picard stricte de mêmes
invariants, laquelle est équivalente à la stricte splittée, et l'équivalence
est en particulier une équivalence de Gr-catégories. On notera que la torsion
ne touche pas $a$, de sorte que la Gr-catégorie sous-jacente n'a pas changé.}

D'où l'énoncé de la page 12, où il faut rétablir deux hypothèses que la page
laisse implicites : \emph{une Gr-catégorie $C$ admet une structure de
catégorie de Picard si et seulement si $\pi_{0}(C)$ est commutatif, agit
trivialement sur $\pi_{1}(C)$, et $k(C) = 0$ ; et elle admet alors aussi une
structure de Picard stricte.}\footnote{La page n'énonce que la condition
$k(P) = 0$, sous la forme « ssi elle est splittable ». Les deux autres sont
nécessaires — elles le sont dès la page 41 du tapuscrit — et sans elles
l'énoncé est faux : une Gr-catégorie de $\pi_{0}$ non commutatif peut être
splittée.}

\pagerange{2}{4}
\subsection*{Multiplier par un groupe, et tordre}

Les pages 2 et 4 fabriquent des catégories de Picard à partir d'une donnée
bilinéaire, et c'est la construction dont le lemme ci-dessus est le moteur.

Soit $C$ une Gr-catégorie et $\Gamma$ un groupe. On regarde $\Gamma$ comme
catégorie discrète d'ensemble sous-jacent $\Gamma$, la loi du groupe
définissant un $\otimes$ d'associativité stricte, et l'on forme $\Gamma \times
C$. On a
\[
\pi_{0}(\Gamma \times C) = \Gamma \times \pi_{0}(C)
\xrightarrow{\ p\ } \pi_{0}(C), \qquad
\pi_{1}(\Gamma \times C) = \pi_{1}(C), \qquad
k_{\Gamma \times C} = p^{*}(k_{C}).
\]
Si $\Gamma$ est commutatif il porte une contrainte de commutativité, et si de
plus $C$ est de Picard, $\Gamma \times C$ l'est.

Donnons-nous alors une application bilinéaire antisymétrique $\varphi : \Gamma
\times \Gamma \to \pi_{1}(C)$. Elle se relève en
$\overline{\varphi}\bigl((\alpha,x),(\beta,y)\bigr) = \varphi(\alpha,\beta)$
sur $\pi_{0}(\Gamma \times C)$, laquelle est encore bilinéaire
antisymétrique, et tord la contrainte de commutativité de $\Gamma \times C$.
La catégorie de Picard obtenue est notée $\Gamma \times^{\varphi} C$ ; elle a
les mêmes $\pi_{0}$ et $\pi_{1}$ que $\Gamma \times C$, et diffère d'elle
exactement par $s$, qui a été augmenté de $\alpha \mapsto
\varphi(\alpha,\alpha)$.

\pagerange{32}{32}
\subsection*{Les droites graduées, ou le plus petit exemple non strict}

La page 32 donne l'exemple, sous le titre « catégorie de Picard-type non
stricte ». Soit $A$ un anneau commutatif unitaire avec $2 \cdot 1_{A} \neq 0$.
Les objets de $C$ sont les $A$-modules gradués modulo $2$ munis d'une base
formée d'un élément, mais donnée seulement au signe près — soit deux éléments
échangés par $x \mapsto -x$, formant chacun une base. Les morphismes sont les
homomorphismes gradués $\varphi : L \to M$ envoyant la base sur la base ; le
produit est le produit tensoriel gradué, avec pour base l'ensemble des $x
\otimes y$ ; la contrainte d'associativité est l'habituelle ; et la contrainte
de commutativité est la \emph{règle de Koszul}.

On a $\pi_{0}(C) = \mathbf{Z}/2\mathbf{Z}$, le degré, et $\pi_{1}(C) =
\{\pm 1\} \cong \mathbf{Z}/2\mathbf{Z}$ — c'est ici que sert l'hypothèse
$2 \cdot 1_{A} \neq 0$, sans laquelle $-1 = 1$ et $\pi_{1}$ serait
nul.\footnote{L'hypothèse est en tête de la page, sans que la page dise à quoi
elle sert.} Et $s$ est l'identité de $\mathbf{Z}/2\mathbf{Z}$ : une droite de
degré impair, échangée avec elle-même, donne $-1$. C'est donc, par le théorème
de la page 42, la plus petite catégorie de Picard non stricte, et à
équivalence près la seule de ces invariants.

La page ajoute une « description alternative » : $C = (\mathbf{Z}/2\mathbf{Z})
\times C_{0}$, le premier facteur catégorie discrète et le second la catégorie
des torseurs sous $\mathbf{Z}/2\mathbf{Z}$. Cette description est exacte comme
description de la \emph{catégorie} et de son $\otimes$, mais elle ne rend pas
la structure de Picard : le produit $\Gamma \times C_{0}$ est strict, et
l'exemple est sa torsion $\Gamma \times^{\varphi} C_{0}$ au sens de la page 4,
avec $\Gamma = \mathbf{Z}/2\mathbf{Z}$ et $\varphi$ l'accouplement de
Koszul.\footnote{Rapprochement qui n'est pas fait sur les pages : la page 32
ne renvoie pas à la page 4, ni l'inverse. Il est proposé ici parce que les
deux constructions coïncident terme à terme et que c'est ce qui explique
pourquoi l'exemple est non strict. La page 32 s'achève sur un paragraphe
« Torseurs gradués » barré de longues diagonales, qui reprend la même
construction sur un topos ; il n'est pas restitué ici.}

Le même exemple avec une graduation par $\mathbf{Z}$ au lieu de
$\mathbf{Z}/2\mathbf{Z}$ a pour invariants $\pi_{0} = \mathbf{Z}$, $\pi_{1} =
\mathbf{Z}/2\mathbf{Z}$ et $s$ la réduction modulo $2$ ; en langue
d'aujourd'hui, c'est le tronqué en degré $1$ du spectre des sphères, et $s$
est la multiplication par $\eta$. On le retrouvera plus bas comme enveloppe de
Picard des $\mathbf{Z}$-modules projectifs de type fini.

\pagerange{50}{52}
\subsection*{La résolution explicite de Sính}

Le manuscrit anglais donne, aux pages 51 et 52, la forme calculatoire de la
classification, et c'est le seul endroit du dossier où elle soit écrite. On y
introduit deux complexes de groupes abéliens libres au-dessus de $M$,
\[
L_{\cdot}(M) : L_{3}(M) \to L_{2}(M) \to L_{1}(M) \to L_{0}(M) \to M,
\]
et de même ${}'L_{\cdot}(M)$, avec
\[
L_{0} = {}'L_{0} = \mathbf{Z}[M], \qquad
L_{1} = {}'L_{1} = \mathbf{Z}[M \times M], \qquad
L_{2} = {}'L_{2} = \mathbf{Z}[M^{3}] + \mathbf{Z}[M^{2}],
\]
\[
{}'L_{3} = \mathbf{Z}[M^{4}] + \mathbf{Z}[M^{3}] + \mathbf{Z}[M^{2}],
\qquad L_{3} = {}'L_{3} + \mathbf{Z}[M],
\]
les différentielles étant celles de la barre, plus $d_{2}[x,y] = [x,y] -
[y,x]$, $d_{3}[x,y] = [x,y] + [y,x]$, et — dans $L_{\cdot}$ seulement —
$d_{3}[x] = [x,x]$. On obtient une bijection canonique entre les classes
d'équivalence de catégories de Picard épinglées de type $(M,N)$ et
$H^{2}\bigl(\mathrm{Hom}({}'L_{\cdot}(M), N)\bigr)$ ; et l'exactitude de
$L_{\cdot}(M)$ donne la trivialité de la classification des strictes.

Les deux complexes ne diffèrent que par le générateur $[x]$ en degré $3$, de
différentielle $[x,x]$, et ce générateur \emph{est} la strictesse : imposer
$c_{X,X} = \mathrm{id}$, c'est annuler la diagonale. D'où la lecture moderne
des deux énoncés. $L_{\cdot}(M)$ est une résolution libre de $M$, donc
$H^{2}(\mathrm{Hom}(L_{\cdot}(M),N)) = \mathrm{Ext}^{2}_{\mathbf{Z}}(M,N)$,
qui est nul parce qu'un groupe abélien a une résolution libre de longueur
$1$ : les catégories de Picard strictes de type $(M,N)$ sont donc toutes
équivalentes, non par accident mais parce que $\mathbf{Z}$ est de dimension
homologique $1$. Et $H^{2}(\mathrm{Hom}({}'L_{\cdot}(M),N)) \cong
\mathrm{Hom}(M, {}_{2}N)$, ce que les notes énoncent directement page 11. On
verra plus bas, page 17, ce que devient la première de ces deux vérités
lorsqu'on quitte le point pour un topos : $\mathrm{Ext}^{2}$ n'y est plus nul,
et il apparaît là où on l'attend.\footnote{Une ligne au crayon en tête de la
page 52, à demi rognée, demande « déterminer $H^{2}(\mathrm{Hom}({}'L_{\cdot}
(M), N))$ » — un lecteur qui pose exactement la question à laquelle la page 11
des notes répond. Ce sont ces complexes que l'on reconnaît aujourd'hui comme
la part de bas degré de la résolution stable d'Eilenberg et Mac Lane.}

\pagerange{7}{10}
\section*{La face géométrique : représentations, 2-gerbes, noyaux}

Les pages 7 à 10 donnent de la classification une seconde lecture, où la
classe $k(C)$ n'est plus le résultat d'un calcul de cocycles mais la classe
d'un objet géométrique. Elles procèdent en trois temps.

\subsection*{Représentations d'un ensemble simplicial}

Soient $K$ un ensemble simplicial et $P$ une $\otimes$-catégorie AU. Une
\emph{représentation} $\Lambda$ de $K$ dans $P$ est la donnée d'un objet
$\Lambda(u) \in P$ pour chaque $u \in K_{1}$ ; pour chaque $\sigma \in K_{2}$,
d'un isomorphisme
\[
\Lambda(d_{2}\sigma) \otimes \Lambda(d_{0}\sigma) \simeq
\Lambda(d_{1}\sigma) ;
\]
et pour chaque $\Delta \in K_{3}$, de la condition de cohérence qui fait
commuter le pentagone associé.\footnote{La page numérote les faces à partir
de $1$ : sa première flèche est $d_{3}\sigma$, sa seconde $d_{1}\sigma$, sa
composée $d_{2}\sigma$. Les indices sont donnés ici dans la convention
d'aujourd'hui, où $d_{2}\sigma$ et $d_{0}\sigma$ sont les deux flèches et
$d_{1}\sigma$ la composée. Avec les lettres de la page — $u = d_{1}\sigma$,
$v = d_{3}\sigma$, $w = d_{2}\sigma$ — l'isomorphisme écrit est
$\Lambda(u) \otimes \Lambda(w) \simeq \Lambda(v)$, qui apparie la composée
avec l'une des deux flèches et ne peut être ce qui est voulu ; c'est
$\Lambda(v) \otimes \Lambda(u) \simeq \Lambda(w)$. La transcription signale
déjà ces indices comme douteux.}

La composée $K_{1} \to P \to \pi_{0}(P) = G$ définit un homomorphisme
$\alpha$, d'où $\overline{\alpha} : K_{1} \to \mathrm{Aut}(\pi_{1}(P))$ et un
système local $\overline{M} = (M, \overline{\alpha})$ sur $K$, où $M =
\pi_{1}(P)$. Le problème est de classer, $K$, $P$ et $\alpha$ étant fixés, les
représentations correspondantes ; la réponse annoncée est qu'elles forment une
2-catégorie qui est un 2-torseur sous une 2-catégorie de Picard stricte
d'invariants
\[
H^{2}(K, \overline{M}), \qquad H^{1}(K, \overline{M}), \qquad
H^{0}(K, \overline{M}),
\]
contravariante en $K$ et covariante en $P$. On la note $\mathrm{Rep}(K; P,
\alpha)$. Les pages 20 à 23, plus bas, démontrent cet énoncé dans le cas qui
sert : celui où $K$ est le nerf d'une catégorie.

\subsection*{La 2-gerbe attachée à une Gr-catégorie}

Soit $P$ une Gr-catégorie d'invariants $G = \pi_{0}(P)$ et $M = \pi_{1}(P)$.
On lui associe un objet sur le topos classifiant $B_{G}$ : à tout $G$-ensemble
$X$, on attache le nerf $\mathcal{X}$ du groupoïde d'action — dont les objets
sont les points de $X$ et les flèches $G \times X$ — muni de
$\alpha_{X} : \mathrm{Fl}\,\mathcal{X} = G \times X \to G$ la projection, et
l'on pose
\[
\widetilde{P}(X) = \mathrm{Rep}(\mathcal{X}; P, \alpha_{X}).
\]
C'est contravariant en $X$, c'est un 2-champ, et c'est une 2-gerbe liée par le
faisceau $M_{B_{G}}$ défini par l'action de $G$ sur $M$.

\textbf{Énoncé} (page 8). \emph{$G$ et le $G$-module $M$ étant fixés, cette
construction est une équivalence entre la 3-catégorie des Gr-catégories
épinglées par $(G,M)$ et la 3-catégorie des 2-gerbes sur $B_{G}$ liées par
$M_{B_{G}}$.}\footnote{L'équivalence est affirmée, non démontrée, sur la page,
et le dossier ne dit pas non plus explicitement que $k(C)$ correspond, sous
elle, à la classe de la 2-gerbe — bien que ce soit la seule lecture cohérente
avec la suite. C'est le dictionnaire que L. Breen a développé et démontré par
la suite, dans ses travaux sur les bitorseurs, la cohomologie non abélienne et
la classification des 2-gerbes et 2-champs.}

Sur les morphismes : si $u : G \to G'$ et $v : M \to M'$ sont compatibles aux
actions, les $\otimes$-homomorphismes $P \to P'$ au-dessus de $(u,v)$ forment
une 2-catégorie, qui est un 2-pseudo-torseur sous une 2-catégorie de Picard
d'invariants $H^{2}(G,M')$, $H^{1}(G,M')$, $H^{0}(G,M')$ ; ce
2-pseudo-torseur est équivalent à celui des homomorphismes de 2-gerbes
correspondants, et l'obstruction à ce qu'il ait un objet est dans
$H^{3}(G,M')$ — c'est la différence des images canoniques des classes de $P$
et de $P'$.

Tout ceci se transporte aux Gr-champs sur un topos $S$ : si $G$ est un groupe
sur $S$ et $M$ un $G$-Module sur $S$, on interprète tout en termes de 2-gerbes
relatives sur $B_{G}$ au-dessus de $S$, c'est-à-dire de 2-gerbes sur $B_{G}$
trivialisées sur $S$ par $S \simeq B_{e} \to B_{G}$. La page note que l'étage
inférieur du même dictionnaire — les 1-gerbes relatives sur $B_{G}$ — répond
aux extensions de $G$ par $M$, et que le cas d'un $M$ non commutatif reste à
traiter.

\subsection*{Les noyaux de Mac Lane}

La page 10 boucle avec la théorie classique. Soit $F$ un groupe et
\[
\varphi : G \longrightarrow \mathrm{Autext}(F)
\]
un homomorphisme — se donner $\varphi$, c'est se donner un \emph{lien} sur
$B_{G}$. On lui associe une Gr-catégorie $P(\varphi)$ d'invariants $G$ et $M =
Z(F)$, comme image inverse par $\varphi$ de la Gr-catégorie $P(F)$ des
bitorseurs sous $F$, dont on a vu que les invariants sont $\mathrm{Autext}(F)$
et $Z(F)$. Les trivialisations de $P(\varphi)$ correspondent alors aux
réalisations de $\varphi$ par une extension de $G$ par $F$ subordonnée à
$\varphi$, et l'on retrouve ainsi l'obstruction de Mac Lane dans $H^{3}(G,M)$.

La page pose ensuite une question qu'il faut lire pour ce qu'elle est : que
toute classe de $H^{3}(G,M)$ s'obtienne à partir d'un noyau ne relève pas,
dit-elle, d'une logique générale. Le fait, lui, est un théorème d'Eilenberg et
Mac Lane ; ce que la page demande n'est pas le fait mais la
raison.\footnote{La formulation de la page est plus ramassée : « le fait que
tous les éléments des $H^{3}(G,M)$ puissent ainsi s'obtenir à l'aide de
"noyaux" ne fait pas (pour l'instant) l'objet d'une logique générale ». La
parenthèse « (pour l'instant) » est de la page. Un paragraphe b), qui devait
donner la description en termes de 2-gerbes, est barré sans avoir été écrit.}

\pagerange{18}{19}
\subsection*{Les splittages, et la Gr-catégorie $C_{0}$}

Reste à dire ce que sont les trivialisations. Soit $C$ une Gr-catégorie
d'invariants $G$, $M$. On note $C_{0}$ la Gr-catégorie splittée de mêmes
invariants :
\[
\mathrm{Ob}\,C_{0} = G, \qquad
\mathrm{Fl}\,C_{0} = G \times M \rightrightarrows G, \qquad
s = b = \mathrm{pr}_{1},
\]
\[
(g,m) \cdot (g,m') = (g, m+m'), \qquad
(g,m) \otimes (g',m') = (gg',\, m + g \cdot m'), \qquad
a_{g,g',g''} = \mathrm{id}.
\]
Un \emph{splittage} de $C$ est une équivalence de Gr-catégories épinglées
$C_{0} \to C$, ce qui revient à la donnée d'un système $L : G \to
\mathrm{Ob}\,C$ et $\varphi : G \times G \to \mathrm{Fl}\,C$ comme au
paragraphe de la classification.\footnote{La page écrit $(g,m) \otimes
(g',m') = (gg', m+m')$, formule qui n'est correcte que si $G$ agit
trivialement sur $M$ ; c'est l'action qui est omise, et elle est rétablie
ici. La page écrit aussi $C$ là où elle définit $C_{0}$, ce que la
transcription signale.}

Les splittages de $C$ forment une catégorie, qui est un 1-torseur sous la
Gr-catégorie $\mathrm{Aut}(C_{0})$ des auto-équivalences épinglées de
$C_{0}$ ; elle est non vide exactement quand $k(C) = 0$. Et les invariants de
$\mathrm{Aut}(C_{0})$ se lisent tout de suite :

\begin{itemize}
\item ses objets sont les 2-cochaînes $\varphi : G \times G \to M$ qui sont
des 2-\emph{cocycles} — la contrainte d'associativité de $C_{0}$ étant
triviale, dire que $(\mathrm{id}, \varphi)$ est un $\otimes$-foncteur, c'est
dire $\delta\varphi = 0$ ;\footnote{La page écrit « qui sont des cobords ».
C'est une erreur, et elle se voit deux lignes plus bas sur la page même : si
les objets étaient les cobords, leur ensemble à isomorphisme près serait
réduit à un point, et non $H^{2}(G,M)$ comme la page l'écrit ensuite. La
transcription note qu'elle porte bien « cobords ».}
\item ses flèches $\varphi \to \psi$ sont les 1-cochaînes de cobord
$\psi - \varphi$ ;
\item d'où $\pi_{0}(\mathrm{Aut}\,C_{0}) = H^{2}(G,M)$ et
$\pi_{1}(\mathrm{Aut}\,C_{0}) = Z^{1}(G,M)$.
\end{itemize}

L'ensemble des splittages à isomorphisme près est donc un torseur sous
$H^{2}(G,M)$, et chacun a $Z^{1}(G,M)$ pour groupe d'automorphismes : c'est,
sous le dictionnaire ci-dessus, la description usuelle de l'espace des
trivialisations d'une 2-gerbe.

La fin de la page 19 fait la même chose au-dessus d'un $G$-ensemble $X$ : on y
construit une 2-catégorie $S(C,X)$ dont les 0-objets sont les systèmes
associant un $L(x,g) \in \mathrm{Ob}\,C$ à chaque couple $(x,g)$ et un
isomorphisme
\[
L(x,g) \otimes L(gx, g') \simeq L(x, g'g)
\]
à chaque triplet. C'est mot pour mot $\mathrm{Rep}(\mathcal{X}; C,
\alpha_{X})$ de la page 8, écrit en clair : la composée de $x
\xrightarrow{g} gx \xrightarrow{g'} g'gx$ est $x \xrightarrow{g'g} g'gx$, et
l'isomorphisme est celui d'une représentation du nerf du groupoïde d'action.

\pagerange{20}{23}
\section*{Extensions d'une catégorie par une Gr-catégorie}

Les pages 20 à 23 sont le calcul qui soutient l'énoncé de la page 7. On y
remplace le groupe par une catégorie quelconque.

\textbf{Définition} (page 20). Soient $\Gamma$ une catégorie et $M$ un groupe.
Une \emph{extension} de $\Gamma$ par $M$ est la donnée, pour toute flèche $u$
de $\Gamma$, d'un $M$-bitorseur $P_{u}$ ; pour tout couple de flèches
composables, d'un isomorphisme $\varphi_{u,v} : P_{u} \otimes P_{v} \simeq
P_{uv}$ ; et de la commutativité, pour trois flèches composables, du
pentagone reliant les deux parenthésages.

Le NB de la page dit ce qui compte : au lieu de $M$ on peut se donner
directement une Gr-catégorie $\widetilde{P}$ — celle des $M$-bitorseurs, par
exemple — et c'est la bonne généralité. Les invariants sont alors un
homomorphisme $\alpha : \mathrm{Fl}(\Gamma) \to \pi_{0}(\widetilde{P})$, qui
vaut $\mathrm{Autext}(M)$ dans le cas des bitorseurs, et une classe $c \in
H^{2}(\Gamma; \widetilde{P})$ — un $H^{2}$ non commutatif, en général dépourvu
d'élément neutre.\footnote{Une marge ajoute : $\mathrm{Autext}(M) =
\mathrm{Aut}(M)$ si $M$ est commutatif. L'absence d'élément neutre est ce qui
distingue ce $H^{2}$ d'un groupe de cohomologie : il n'y a pas d'extension
triviale privilégiée, ce qui est précisément la situation de Mac Lane.}

Une \emph{2-extension} remplace les bitorseurs $P_{u}$ par des $M$-gerbes, les
isomorphismes $\varphi_{u,v}$ par des équivalences, et la commutativité du
pentagone par un isomorphisme $\alpha_{u,v,w}$ entre les deux composés. Il
faut alors, et alors seulement, un axiome de plus, pour quatre flèches
composables : que tous les isomorphismes
\[
\bigl(\cdots((P_{u_{1}} \otimes P_{u_{2}}) \otimes P_{u_{3}}) \cdots\bigr)
\longrightarrow P_{u_{1}u_{2}\cdots u_{n}}
\]
déduits des $\varphi$ et de l'associativité forment un système
transitif.\footnote{La page 22 présente cette condition comme un axiome d) sans
distinguer les deux cas, et c'est la distinction qui manque. Pour une
extension au sens de la page 20, où le pentagone \emph{commute}, la
transitivité pour $n$ facteurs est une conséquence du cas $n = 3$ par
l'argument de cohérence de Mac Lane, et non un axiome. Pour une 2-extension,
où le pentagone ne commute qu'à un $\alpha_{u,v,w}$ près, c'en est un
véritable : c'est la condition de 3-cocycle sur les $\alpha$. La marge de la
page porte une colonne biffée d'essais de parenthésages, qui est le calcul
correspondant.}

\subsection*{La 2-catégorie des $\alpha$-extensions}

Soient $\Gamma$ une catégorie, $P$ une Gr-catégorie, et $\alpha :
\mathrm{Fl}(\Gamma) \to \pi_{0}(P)$ un homomorphisme, de composé
$\overline{\alpha}$ dans $\mathrm{Aut}(\pi_{1}(P))$. La page 23 décrit la
2-catégorie des extensions de $\Gamma$ par $P$ subordonnées à $\alpha$ :

\begin{itemize}
\item ses 0-objets sont ces extensions ;
\item ses 1-objets $\{P_{u}, \varphi_{u,v}\} \to \{P'_{u},
\varphi'_{u,v}\}$ sont les systèmes d'homomorphismes $\lambda(u) : P_{u} \to
P'_{u}$ faisant commuter le carré
\end{itemize}

\begin{tikzcd}[column sep=normal, row sep=small, nodes={font=\scriptsize}]
P_u \otimes P_v \arrow[r, "\varphi_{u{,}v}"]
\arrow[d, "\lambda(u) \otimes \lambda(v)"'] & P_{uv} \arrow[d, "\lambda(uv)"] \\
P'_u \otimes P'_v \arrow[r, "\varphi'_{u{,}v}"'] & P'_{uv}
\end{tikzcd}

\begin{itemize}
\item et ses 2-objets, de $\lambda$ à $\mu$, sont les 0-cochaînes $m$ sur
$\Gamma$ à valeurs dans $\pi_{1}(P)$ dont le cobord réalise le passage de
l'un à l'autre.
\end{itemize}

Le calcul qui donne l'énoncé est celui-ci. Si $\{\lambda(u)\}$ est un
1-objet, tous les autres sont de la forme $\lambda(u) \otimes f(u)$ avec
$f : \mathrm{Fl}\,\Gamma \to \pi_{1}(P)$ vérifiant
\[
f(uv) = f(u) \cdot \bigl(\alpha(u) \cdot f(v)\bigr),
\]
c'est-à-dire $f \in Z^{1}\bigl(\Gamma, (\pi_{1}(P),
\overline{\alpha})\bigr)$ ; l'ensemble des 1-objets est donc un
pseudo-torseur sous $Z^{1}$. Et $\mu = \lambda \otimes f$ est reliée à
$\lambda$ par les $m$ tels que $f = \delta m$, où
$(\delta m)(u) = m(y) - \alpha(u) \cdot m(x)$ pour $u : x \to
y$.\footnote{La page écrit $m \in Z^{0}(\Gamma, \pi_{1}(P))$ ; il s'agit des
0-cochaînes, c'est-à-dire des fonctions sur les objets de $\Gamma$, dont les
cocycles proprement dits sont le noyau de $\delta$. La page achève la formule
sur un point d'interrogation, qui est d'elle.}

En passant aux classes d'isomorphisme, on obtient exactement ce que la page 7
annonçait : la 2-catégorie des $\alpha$-extensions est un 2-pseudo-torseur
d'invariants $H^{2}$, $H^{1}$, $H^{0}$ de $\Gamma$ à valeurs dans le système
local $(\pi_{1}(P), \overline{\alpha})$. Une extension de $\Gamma$ par $P$
subordonnée à $\alpha$ est une représentation du nerf de $\Gamma$ dans $P$ ;
les deux énoncés sont le même.\footnote{La page 21 est une feuille de calcul
préparatoire sans prose, couverte de colonnes simpliciales, d'accouplements
$P \otimes Q \to R$ sous leurs trois ou quatre formes, et d'un tétraèdre dont
la transcription note qu'elle ne peut orienter les arêtes. Rien n'en est
restitué ici : ce qu'elle établit est ce qui précède, et une feuille d'essais
dont plusieurs indices sont donnés comme douteux ne peut être reproduite dans
une édition de lecture sans y faire figurer un calcul que personne ne peut
vérifier.}

\pagerange{13}{17}
\section*{Le passage au topos : les champs de Picard}

Les pages 13 à 17 refont toute la théorie commutative au-dessus d'un topos
$S$, où $\pi_{0}$ et $\pi_{1}$ deviennent des faisceaux abéliens $M$, $N$. Le
cadre est celui de Deligne : les champs de Picard \emph{stricts} sur $S$
correspondent aux complexes de la catégorie dérivée $D(\mathrm{Ab}(S))$
concentrés en deux degrés, et leurs $\mathrm{Hom}$ aux $R\mathrm{Hom}$
tronqués. Ce que le dossier ajoute est le cas non strict, où le triplet
$(M, N, \sigma : M \to {}_{2}N)$ remplace le couple.

\textbf{Les foncteurs.} Soient $P$, $P'$ des champs de Picard épinglés par
$(M,N,\sigma)$ et $(M',N',\sigma')$, et un homomorphisme de triplets fixé.
Le champ des $\otimes$-foncteurs AUC $P \to P'$ qui l'induisent est un
1-torseur sous le champ de Picard strict d'invariants $\mathrm{Ext}^{1}(M,N')$
et $\mathrm{Hom}(M,N')$, obtenu par troncature de $R\mathrm{Hom}(M,N')$.
Trois conséquences, que la page tire successivement :

\begin{itemize}
\item deux foncteurs induisant le même homomorphisme de triplets ne sont pas
nécessairement isomorphes ; ils le sont si $\mathrm{Ext}^{1}(S; M,N') = 0$, et
localement isomorphes si le faisceau $\underline{\mathrm{Ext}}^{1}(M,N')$ est
nul ;
\item tout homomorphisme de triplets n'est pas nécessairement induit ;
l'obstruction est dans $\mathrm{Ext}^{2}(S; M,N')$, elle s'annule localement,
et vit donc dans $H^{0}(S, \underline{\mathrm{Ext}}^{2}(M,N'))$ ;
\item et pour $(M,N,\sigma)$ fixé, les champs de Picard de ces invariants
forment une 2-catégorie qui est un 2-pseudo-torseur sous celle des champs de
Picard \emph{stricts} d'invariants $(M,N)$ : la « différence » de deux d'entre
eux est un champ de Picard strict d'invariants $M$, $N$, c'est-à-dire une
classe de $\mathrm{Ext}^{2}(S; M,N)$.
\end{itemize}

Ce dernier point est exactement le dictionnaire de Deligne lu à l'envers : un
complexe concentré en deux degrés, de cohomologies $N$ et $M$, s'inscrit dans
un triangle distingué $N[1] \to K \to M \to N[2]$, et de tels complexes sont
classés par $\mathrm{Hom}_{D}(M, N[2]) = \mathrm{Ext}^{2}(M,N)$.

\textbf{L'existence.} Reste à savoir quels triplets $(M,N,\sigma)$ proviennent
effectivement d'un champ de Picard, et c'est la question à laquelle le lemme
de la page 4 répondait sur le point. Elle est claire si $\sigma$ est la
restriction d'une application bilinéaire antisymétrique $M \times M \to N$,
donc si $M/2M$ est un faisceau « libre », par exemple si $M$ est libre sur
$\mathbf{Z}$. En général, choisissons une présentation $0 \to R \to L
\xrightarrow{p} M \to 0$ avec $L$ tel que $\sigma \circ p$ soit réalisable par
un champ $P$, muni d'une classe de $\mathrm{Ext}^{2}(S; L,N)$. La restriction
de $P$ à $R$ définit une classe de $\mathrm{Ext}^{2}(S; R,N)$, déterminée
modulo les images des éléments de $\mathrm{Ext}^{2}(S; L,N)$, dont l'image
\[
\omega(\sigma) \in \mathrm{Ext}^{3}(S; M,N)
\]
est canoniquement déterminée ; et $\omega(\sigma) = 0$ si et seulement si l'on
peut choisir $P$ tel que $P|_{R}$ soit splitté, c'est-à-dire tel que $P$
provienne d'un champ d'invariants $M$, $N$. On vérifie que $\omega(\sigma)$ ne
dépend pas de la présentation choisie, qu'elle est fonctorielle en $M$ et $N$,
et compatible aux images inverses par morphismes de topos.

\textbf{Et l'obstruction est nulle.} La page 16 conclut par l'argument
d'universalité : on se place sur le topos classifiant $B$ de la donnée
$\sigma : M \to {}_{2}N$, où $\mathrm{Ext}^{i}(B; M,N) = 0$ pour $i > 0$ ; en
particulier $\omega(\sigma) \in \mathrm{Ext}^{3}(B;M,N) = 0$, et la
fonctorialité transporte cette nullité partout. Tout triplet $(M,N,\sigma)$
provient donc d'un champ de Picard.\footnote{L'annulation des
$\mathrm{Ext}^{i}$ sur le topos classifiant est affirmée sur la page — le mot
qui la porte, « aurions », est donné comme lecture incertaine par la
transcription — et n'y est pas établie. La conclusion, elle, est corroborée
par la page 17 : la suite exacte qu'on y lit est surjective sur
$\mathrm{Hom}(M, {}_{2}N)$, ce qui dit la même chose. En marge de la page 15,
une glose : cette condition s'interprète en disant que $P$ admet une
contrainte de commutativité qui en fasse un champ de Picard \emph{strict} ;
la page ne dit pas de quel $P$ elle parle, et on ne le restitue pas ici.}

\textbf{Le bilan de la page 17.} Soit $\mathcal{P}$ la 2-catégorie des champs
de Picard sur $X$ épinglés par $(M,N)$, munie de la composition de Baer. Ses
invariants sont
\[
0 \longrightarrow \mathrm{Ext}^{2}(M,N) \longrightarrow \pi_{0}(\mathcal{P})
\longrightarrow \mathrm{Hom}(M, {}_{2}N) \longrightarrow 0,
\]
\[
\pi_{1}(\mathcal{P}) = \mathrm{Ext}^{1}(M,N), \qquad
\pi_{2}(\mathcal{P}) = \mathrm{Hom}(M,N),
\]
où $\pi_{0}$ est le groupe des objets à équivalence près, $\pi_{1}$ celui des
classes d'isomorphie d'auto-équivalences d'un objet, et $\pi_{2}$ celui des
automorphismes du foncteur identique d'un objet.\footnote{La page écrit
$\pi_{0}$ deux fois, la seconde pour le groupe qui est en fait le $\pi_{2}$ ;
la transcription le signale. Elle s'achève sur une esquisse fragmentaire —
« $\mathcal{P}$ 2-catégorie de Picard stricte ? », suivie d'un début de
description de $P \otimes P$ — dont plusieurs symboles sont illisibles ; elle
n'est pas restituée.}

C'est le point où toutes les pages du dossier se recoupent. La suite exacte
dit que $\pi_{0}(\mathcal{P})$ est une extension de la classification par
$\sigma$ — celle de Sính — par le groupe $\mathrm{Ext}^{2}(M,N)$ qui mesure
les champs stricts. Sur le topos ponctuel, $\mathrm{Ext}^{2}_{\mathbf{Z}}(M,N)
= 0$ et la suite donne $\pi_{0}(\mathcal{P}) \cong \mathrm{Hom}(M, {}_{2}N)$ :
on retrouve exactement l'énoncé des pages 11 et 42. Et $\pi_{1} =
\mathrm{Ext}^{1}(M,N)$, $\pi_{2} = \mathrm{Hom}(M,N)$ sont les deux invariants
que la page 11 attribuait déjà à la catégorie de Picard sous laquelle les
équivalences forment un torseur. Ce qui, sur le point, était un théorème sans
reste devient sur un topos une suite exacte à deux termes, et le terme
nouveau est exactement l'obstruction homologique que le point n'avait pas.

\pagerange{31}{37}
\section*{Suites de fibrations et tours de lacets}

Trois feuilles du dossier — les pages 31, 35 et 37 — sont des feuilles de
calcul sans prose, couvertes de treillis de suites de fibrations. Elles ne
sont pas restituées ici sous leur forme graphique, et c'est une décision qu'il
vaut mieux énoncer que taire.\footnote{Les treillis comportent une dizaine de
colonnes en $\Omega^{k}$ répétées deux ou trois fois, et des escaliers de
$X_{ij}$ dont la transcription donne plusieurs cases comme surchargées ou
douteuses — un $X_{10}$ qui se lit aussi bien $X_{01}$, une case entièrement
récrite. La transcription les rend structurellement, case par case, avec ces
réserves ; c'est là qu'il faut aller les voir. Les reproduire ici les
présenterait comme des diagrammes établis, ce qu'ils ne sont pas, et une
édition de lecture qui affiche un treillis invérifiable ment sur ce qu'elle
peut promettre. Ce qui suit est ce que ces feuilles \emph{disent}, et rien de
plus.}

Ce qu'elles disent est ceci. Une fibration $F \to X \to S$ dans un cadre
pointé engendre sa suite de Puppe,
\[
\cdots \to \Omega^{2}S \to \Omega F \to \Omega X \to \Omega S \to F \to X
\to S,
\]
et les pages la font proliférer dans deux directions à la fois : chaque rangée
d'un treillis est une telle suite, chaque colonne en est une autre, et chaque
bloc de rangées reprend le précédent après application de $\Omega$. La page 37
fait la même chose pour un carré $X_{ij}$ : les suites de fibrations des
lignes et des colonnes s'engrènent, chaque tour de trois rangées reprenant le
précédent sous $\Omega$. C'est le calcul qui organise, en langue
d'aujourd'hui, la suite spectrale d'une tour de fibrations.

La page 35 en donne la version attachée à des flèches composables $A
\xrightarrow{f} B \xrightarrow{g} C \xrightarrow{h} D$ : les fibres
$\Gamma_{f}$, $\Gamma_{gf}$, $\Gamma_{hgf}$, \ldots{} y sont reliées par des
suites qui se déduisent l'une de l'autre par $\Omega$ d'un côté et par
$\Sigma$ de l'autre. C'est l'axiome de l'octaèdre sous sa forme itérée : la
fibre d'une composée s'inscrit dans une suite de fibrations avec les deux
fibres.

La même page énonce en clair, au centre, le critère pour qu'un carré soit
cartésien au sens homotopique. Étant donné $Z$ au-dessus de $S$ avec ses deux
projections $p$, $q$ vers $X$ et $Y$, sont équivalentes les conditions :
$\Gamma_{g} \to \Gamma_{f}$ est un isomorphisme ; $\Gamma_{p} \to \Gamma_{g}$
en est un ; et la propriété universelle faible — pour tout $T$ muni de $p' : T
\to X$, $q' : T \to Y$ avec $fp' = gq'$, il existe $u : T \to Z$ avec $p' =
pu$ et $q' = qu$, sans que $u$ soit nécessairement
unique.\footnote{L'intitulé porté sur la page est donné par la transcription
comme « Conditions équivalentes à "$S$ connexe" », lecture incertaine et qui
ne correspond pas aux conditions énoncées : celles-ci sont les critères
usuels pour que le carré soit cartésien. Une quatrième condition, portant sur
$\Gamma_{r}$, est trop raturée pour être restituée. La non-unicité de $u$ est
soulignée par la page elle-même, et c'est ce qui distingue une limite
homotopique d'une limite.}

Enfin, cerclée au bas de la page 35, une remarque qui vaut plus que le calcul
qui l'entoure : se donner le foncteur $\Omega$ et une notion de \emph{suite
distinguée} $\Omega S \to F \to X \to S$, avec ses axiomes. C'est le programme
d'une axiomatisation du cadre où toutes ces suites vivent — la moitié
« fibres » de ce qui deviendra la structure des catégories stables et, dans le
vocabulaire plus tardif de Grothendieck lui-même, des dérivateurs. La page
n'écrit pas les axiomes.

\pagerange{25}{29}
\section*{Le programme des pages 25 à 29}

Les pages 25 à 29 sont d'une autre nature que tout le reste : une liste
numérotée de onze points, d'une écriture très abrégée, qui est un programme de
travail. On la restitue dans son ordre.

\begin{enumerate}
\item Les $n$-catégories, strictes ou non, et les $n$-foncteurs ; elles
forment une $(n{+}1)$-catégorie dont les $\mathrm{Hom}$ sont des
$n$-catégories ; notion de $i$-fidélité, la $(n{+}1)$-fidélité s'appelant
$n$-équivalence. Une Gr-catégorie « est » une 2-catégorie à un seul 0-objet.
Il y a des chances, dit la page, que les non strictes se ramènent aux
strictes.
\item Les exemples venus des structures simpliciales, et le cœur du point :
pour $A$ une catégorie et $C$ une catégorie, se donner un $\mathrm{Hom}$
interne bifonctoriel à valeurs dans $\hat{C}$ équivaut à se donner un foncteur
$\hat{C} \to \mathrm{Hom}_{!}(\hat{A}, \hat{A})$ commutant aux limites
inductives, la composition des $\mathrm{Hom}$ correspondant à la
multiplicativité de ce foncteur et l'associativité à sa
compatibilité.\footnote{Autrement dit : une catégorie enrichie sur $\hat{C}$
est la même chose qu'une action du produit de convolution de $\hat{C}$ sur
$\hat{A}$. La page passe par trois écritures successives de la même donnée,
que la transcription rend telles quelles ; c'est celle-ci qu'elles disent
toutes.} Pour $C$ la catégorie des ensembles finis totalement ordonnés non
vides, $\hat{C}$ est celle des ensembles simpliciaux, on dispose dans les
$\mathrm{Hom}$ de notions d'homotopie, et l'on associe à $A$ des
$n$-catégories pour tout $n$ — construction facile si les $\mathrm{Hom}$ sont
de Kan.\footnote{Le dossier écrit « ensembles $\tfrac12$-simpliciaux » pour
ce qu'on appelle aujourd'hui ensembles simpliciaux. C'est le chemin qui mène
aux catégories simplicialement enrichies, et de là aux $(\infty,1)$-catégories
telles qu'on les pratique.} La page ajoute, sous un numéro 2$'$, les topos de
Quillen et l'interprétation des $n$-catégories strictes comme structures
simpliciales particulières.
\item Les $n$-catégories fibrées et les $n$-champs sur un site ou un
topos : équivalence entre $n$-champs sur un site et sur le topos associé ;
$n$-champ enveloppant d'une $n$-catégorie fibrée ; existence des images
inverses, déduites des images directes par la formule d'adjonction.
\item Les $n$-champs de profondeur $\geq i$ le long d'un fermé.
\item Les formules de changement de base, d'une topologie ordinaire à une
topologie étale.
\item Les théorèmes du type de Lefschetz, en topologie étale ou dans
l'analytique complexe.
\item Les $n$-champs localement constants : un morphisme de topos $f : X
\to Y$ est une équivalence d'homotopie si et seulement si il induit une
équivalence $\Gamma(Y,A) \simeq \Gamma(X, f^{*}A)$ pour tout $n$-champ
localement constant $A$ — d'où l'idée qu'un type d'homotopie de topos est
reconnu par la donnée, pour tout $n$, de la $(n{+}1)$-catégorie des
$n$-champs localement constants, avec les foncteurs de passage entre étages ;
dans le cadre d'Artin et Mazur, avec les restrictions de constructibilité
voulues.
\item Les structures algébriques sur les $n$-catégories et $n$-champs :
Gr-structures, structures de Picard, de Picard stricte, torseurs, et leur
stabilité par images directes et inverses.
\item La $(n{+}1)$-catégorie des $n$-champs de Picard \emph{stricts} sur
$X$ est équivalente à la sous-catégorie pleine de $D(\mathrm{Ab}(X))$ formée
des complexes concentrés en $n{+}1$ degrés consécutifs ; les opérations,
images directes et inverses se correspondent, à une troncature près pour
l'image directe, et les $\mathrm{Hom}$ correspondent aux $R\mathrm{Hom}$
tronqués.\footnote{La page écrit la condition $\underline{H}^{i}(K^{\cdot}) =
0$ pour $i \notin [0,n]$, en indexation cohomologique, alors que le tapuscrit
de la page 42 indexe homologiquement ($H_{0} = \pi_{0}$, $H_{1} = \pi_{1}$).
L'un est l'autre au signe des degrés près, et l'on ne fixe pas ici de
convention pour un énoncé programmatique. Le cas $n = 1$ est le théorème de
Deligne ; le cas général est ce qu'on appelle aujourd'hui la correspondance de
Dold-Kan stable, et les objets de droite sont les spectres connectifs
$n$-tronqués en modules sur $\mathbf{Z}$.} La page ajoute qu'il y aurait lieu
d'interpréter géométriquement les objets tronqués de $D(\mathrm{Mod}(O))$ pour
$O$ un Anneau sur $X$, ainsi que $Rf_{*}$, $Lf^{*}$ et $R\mathrm{Hom}$, et de
lire les $\otimes^{L}$ comme des accouplements universels de catégories de
Picard.
\item L'étude des $n$-Gr-champs. Le cas $n = 1$ est fait, pour $X$
quelconque ; pour $n = 2$ et $X$ le topos ponctuel, la page ne sait quoi
dire.\footnote{« cf lettre à Deligne », renvoie la page — lecture donnée comme
incertaine. La correspondance avec Deligne et avec Breen est nommée deux fois
dans le dossier, ici et page 29.}
\item Le $n$-champ de Picard \emph{enveloppant} d'un $n$-$\otimes$-champ
ACU, dont les invariants $\pi_{i}$ seraient les faisceaux $K_{i}(A)$.
\end{enumerate}

Et la page 29 pose la question ouverte : que dire de la structure des
$n$-champs de Picard \emph{non} stricts ? Le cas $n = 1$ se voit — c'est tout
ce qui précède — mais rien pour $n \geq 2$.

Deux de ces points ont trouvé leur réponse, et il vaut de dire laquelle. Le
point 11 est le théorème de complétion en groupe : la construction
enveloppante d'une catégorie monoïdale symétrique est son spectre de
K-théorie, et ses $\pi_{i}$ sont bien les $K_{i}$ — ce que Segal et May
établissent, par les $\Gamma$-espaces et par les opérades $E_{\infty}$, dans
les années mêmes du dossier. Et la question de la page 29 se reformule ainsi :
les $n$-champs de Picard non stricts sont les spectres connectifs $n$-tronqués,
dont la structure est celle de leur tour de Postnikov, avec des invariants de
Postnikov stables ; le cas $n = 1$ n'a qu'un seul invariant, et c'est le $s$
des pages 12 et 42, autrement dit la multiplication par $\eta$. C'est
précisément parce que $s$ ne s'annule pas que le point 9 doit se restreindre
aux champs \emph{stricts}, qui sont les modules sur $\mathbf{Z}$ — et donc les
complexes.

\pagerange{42}{44}
\section*{L'enveloppe de Picard : $K^{0}$ et $K^{1}$}

Le dernier tiers du tapuscrit résout un problème universel, et c'est là que le
dossier rejoint la K-théorie.

Soient $C$ une $\otimes$-catégorie AUC et $S$ une partie de
$\mathrm{Ob}\,C$. On cherche à envoyer $C$, par un $\otimes$-foncteur AUC,
dans une $\otimes$-catégorie AUC $G$ de manière que $F(X)$ soit inversible
pour tout $X \in S$, et universellement. Le résultat principal est que ce
problème a toujours une solution, notée $S^{-1}C$, et que pour toute $G$ le
foncteur naturel
\[
\mathrm{Hom}_{\otimes \mathrm{AUC}}(S^{-1}C, G) \longrightarrow
\mathrm{Hom}^{S}_{\otimes \mathrm{AUC}}(C, G)
\]
est un isomorphisme de catégories, l'exposant $S$ désignant la sous-catégorie
pleine des foncteurs qui inversent les objets de $S$.

Lorsque $S$ est l'ensemble de \emph{tous} les objets de $C$, la catégorie
$S^{-1}C$ est elle-même une catégorie de Picard, non nécessairement stricte :
c'est l'\emph{enveloppe de Picard} de $C$, et ses invariants méritent, dit le
tapuscrit, la notation
\[
K^{0}(C) = \pi_{0}(S^{-1}C), \qquad K^{1}(C) = \pi_{1}(S^{-1}C).
\]

\textbf{Théorème.} \emph{Si $C$ est la catégorie des modules projectifs de
type fini sur un anneau $R$, munie de la somme, ces groupes sont les
invariants classiques $K^{0}(R)$ et $K^{1}(R)$ — le groupe de Grothendieck et
le groupe de Whitehead.}\footnote{Page 43 pour le tapuscrit, page 56 pour le
manuscrit anglais, lequel écrit « the category of all finitely generated
$R$-modules » en omettant \emph{projectifs} ; c'est de modules projectifs de
type fini qu'il s'agit, comme le tapuscrit l'écrit correctement, faute de quoi
l'énoncé est faux. Sur la page 43, les mots « on peut démontrer que » ont été
barrés d'un trait ondulé, et le tapuscrit ajoute que le même résultat vaut
sans doute pour toute catégorie additive munie de la somme.}

Vient alors la remarque la plus fine du tapuscrit, et elle est négative. Cette
catégorie stabilisée n'est pratiquement jamais une catégorie de Picard
stricte, même dans le cas des modules projectifs de type fini ; or, par le
dictionnaire de Deligne, une catégorie de Picard stricte est un complexe. Il
est donc peu probable qu'on puisse construire dans la catégorie dérivée de
$(\mathrm{Ab})$ un objet canonique $K_{\cdot}$ avec $H_{0}(K_{\cdot}) = K^{0}$
et $H_{1}(K_{\cdot}) = K^{1}$ : une telle construction fournirait précisément
la catégorie de Picard stricte qui n'existe pas.

En langue d'aujourd'hui, cet argument est exact et il nomme l'obstruction.
Que l'enveloppe ne soit pas stricte, c'est que $s \neq 0$ ; et $s$ est la
multiplication par $\eta$, de $K^{0}$ dans $K^{1}$. Un complexe de chaînes
donnerait un module sur $\mathbf{Z}$, où $\eta$ agit par zéro. Le spectre de
K-théorie n'en est pas un, et c'est pourquoi il ne se laisse pas remplacer par
un complexe.\footnote{On peut le voir sur le plus petit exemple. Pour $R =
\mathbf{Z}$ on a $K^{0} = \mathbf{Z}$, $K^{1} = \{\pm 1\} \cong
\mathbf{Z}/2\mathbf{Z}$, et $s$ est la réduction modulo $2$. Par le théorème
de la page 42, l'enveloppe de Picard des $\mathbf{Z}$-modules projectifs de
type fini est donc équivalente à la catégorie des droites graduées de la page
32, dans sa version graduée par $\mathbf{Z}$ — l'objet de degré $n$ étant le
module de rang $n$, et le signe étant celui que la permutation des facteurs
introduit sur les bases. Les deux pages ne se citent pas.}

Le tapuscrit termine sur la question suivante, qui est le point 11 du
programme : une interprétation géométrique analogue pour les $K^{i}$
supérieurs, à chercher du côté des $n$-catégories de Picard et de leurs
enveloppes.

\pagerange{52}{56}
\subsection*{Les deux problèmes universels du chapitre III}

Le manuscrit anglais donne la construction que le tapuscrit se dispense
d'expliciter — « fort longue et fastidieuse », dit-il. Elle passe par deux
problèmes universels dont le second se réduit au premier.

\textbf{Rendre des objets unités.} Soient $A$ une $\otimes$-catégorie AC, $A'$
une $\otimes$-catégorie AC dont la catégorie de base est un groupoïde, et
$(T,\check{T}) : A' \to A$ un $\otimes$-foncteur AC. On cherche une
$\otimes$-catégorie ACU $P$, un $\otimes$-foncteur AC $(D,\check{D}) : A \to
P$, et un $\otimes$-isomorphisme $\lambda$ de $D \circ T$ sur le foncteur
constant $1_{P}$, universels.

La construction prend $\mathrm{Ob}\,P = \mathrm{Ob}\,A$ et, pour morphismes,
\[
\mathrm{Hom}_{P}(A,B) = \Phi(A,B)/R_{A,B},
\]
où $\Phi(A,B)$ est l'ensemble des triplets $(A', B', u)$ avec $A', B' \in
\mathrm{Ob}\,A'$ et $u : A \otimes TA' \to B \otimes TB'$ dans $A$, et où
$R_{A,B}$ identifie deux triplets lorsqu'il existe $C'_{1}$, $C'_{2}$ et des
isomorphismes $A'_{1} \otimes C'_{1} \simeq A'_{2} \otimes C'_{2}$, $B'_{1}
\otimes C'_{1} \simeq B'_{2} \otimes C'_{2}$ rendant commutatif, \emph{dans la
catégorie quotient $A^{S}$}, le diagramme de comparaison.\footnote{$A^{S}$ est
le quotient de $A$ par la partie multiplicative $S$ engendrée par les
endomorphismes $T(c_{A',A'})$ — les symétries des carrés. C'est là le nerf de
la construction : ce sont exactement les signes que le passage au quotient
efface, et un lecteur l'a bien vu, qui a porté au crayon en marge de la page
52 « what about composition ? » et corrigé « endomorphisms » en « arrows ».}
Le manuscrit donne ensuite la composition, la structure $\otimes$, les
contraintes ACU, le foncteur $D$ et l'isomorphisme $\lambda$, chacun par un
diagramme explicite ; on ne les reproduit pas ici, la transcription les
portant en entier.

\textbf{Inverser des objets.} Soient $C$, $C'$ des $\otimes$-catégories ACU,
$C'$ de base un groupoïde, et $(F,\check{F}) : C' \to C$ un
$\otimes$-foncteur ACU. On cherche $P$ et $(Q,\check{Q}) : C \to P$ tels que
$QFX'$ soit inversible pour tout $X'$, universellement. Le problème se ramène
au premier en posant $A' = C'$, $A = C \times C'$ et $TX' = (FX', X')$,
grâce à l'équivalence canonique
\[
\mathrm{Hom}^{\otimes \mathrm{ACU}}(C \times C', Q) \longrightarrow
\mathrm{Hom}^{\otimes \mathrm{ACU}}(C, Q) \times
\mathrm{Hom}^{\otimes \mathrm{ACU}}(C', Q).
\]
La catégorie $P$ obtenue est la \emph{$\otimes$-catégorie de fractions} de
$C$ relativement à $(C', F)$. Prenant pour $C'$ le sous-groupoïde maximal
$C^{is}$ et pour $F$ l'inclusion, on obtient l'enveloppe de Picard
$\mathrm{Pic}(C)$, et
\[
\pi_{0}\bigl(\mathrm{Pic}(P(R))\bigr) \simeq K^{0}(R), \qquad
\pi_{1}\bigl(\mathrm{Pic}(P(R))\bigr) \simeq K^{1}(R).
\]
En marge de cette formule, de la main de l'annotateur, une phrase en anglais
qui dit exactement ce que le tapuscrit disait à la page 43 : cette enveloppe
de Picard n'est pas une catégorie de Picard stricte.\footnote{Deux lecteurs,
deux notations, une seule remarque — c'est le genre de recoupement que seule
la lecture du dossier entier fait apparaître. La construction elle-même est de
la même famille que celles qui portent aujourd'hui le nom de complétion en
groupe : le $S^{-1}S$ de Quillen, ou $\Omega B$ appliqué au classifiant.}

\pagerange{57}{57}
\subsection*{La catégorie de suspension}

La dernière page mathématique du manuscrit anglais tire de tout ceci une
conséquence négative, et c'est la seule du dossier qui porte sur la topologie.

Soient $C$ une $\otimes$-catégorie ACU, $Z$ un objet distinct de l'unité, et
$S = -\otimes Z$. La \emph{catégorie de suspension} de $C$ définie par $Z$ est
la solution $(P, i, p)$ du problème universel des triples $(Q, j, q)$ où $q$
est une auto-équivalence de $Q$ et où le carré reliant $S$ et $q$ commute à
isomorphisme naturel près. Pour $C$ la catégorie homotopique des espaces
pointés munie du smash et $Z = S^{1}$, on retrouve la définition classique de
la catégorie de suspension.

Soient alors $C'$ la sous-catégorie $\otimes$-stable engendrée par $Z$ et $P$
la $\otimes$-catégorie de fractions de $C$ relativement à $C'$. On dispose
d'un foncteur $G$ de la catégorie de suspension vers $P$. Si $G$ n'est pas
fidèle — ce qui, affirme la page, est le cas pour les espaces pointés et
$Z = S^{1}$ — alors il est impossible de munir la catégorie de suspension
d'une loi $\otimes$ qui en fasse une $\otimes$-catégorie ACU avec $iZ$
inversible et $i$ un $\otimes$-foncteur ACU.\footnote{L'affirmation de non
fidélité n'est pas démontrée sur la page, et elle ne l'est pas davantage ici.
Ce que la page met en évidence est la différence entre inverser un objet et
l'inverser \emph{symétriquement} : c'est le point sur lequel butent les
constructions naïves d'une catégorie monoïdale symétrique de spectres, et ce
qui a fini par imposer les constructions modernes plutôt qu'une simple
inversion de $S^{1}$.}

\pagerange{44}{44}
\section*{Une algèbre homologique non commutative}

Le dossier finit sur une prophétie, et elle n'a besoin d'aucun préalable. Les
développements esquissés, dit la page 44, obligeront à tirer au clair les
relations très étroites qu'on pressent entre trois séries de notions : les
$n$-catégories et $\infty$-catégories d'un côté, les ensembles simpliciaux et
les espaces topologiques de l'autre ; de même entre les $n$-Gr-catégories et
les groupes simpliciaux ou topologiques ; de même enfin entre les
$n$-catégories de Picard \emph{strictes} et les complexes de chaînes tronqués
à la dimension $n$. Il y aurait à s'attendre à une fusion plus ou moins
complète entre ces trois visions — algèbre homologique, algèbre homotopique,
algèbre catégorique — dans une vision commune, qui pourrait bien prendre le
nom d'algèbre homologique non commutative,

\begin{quote}
étant entendu que ce qui porte actuellement ce nom n'est qu'une amorce très
partielle de ce qui doit venir\ldots
\end{quote}

Suit une dernière question, « plus terre à terre ». Si $R$ est commutatif, ou
plus généralement si $C$ est une catégorie additive munie d'un $\otimes$ ACU
biadditif, l'enveloppe de Picard porte, outre sa somme $\oplus$, une seconde
opération $\otimes$ venue du produit tensoriel, reliée à la première par des
isomorphismes de distributivité
\[
X \otimes (Y \oplus Z) \simeq (X \otimes Y) \oplus (X \otimes Z).
\]
Il y aurait lieu d'étudier systématiquement les catégories ainsi munies de
deux opérations AUC et d'une donnée de distributivité, et d'en faire une
théorie de structure lorsqu'elles sont de Picard pour $\oplus$. Cela implique
en particulier que $\pi_{0}$ est un anneau commutatif unitaire et $\pi_{1}$ un
module sur cet anneau ; et la question posée est : \emph{quels autres
invariants faut-il introduire pour caractériser une telle catégorie à
équivalence près, respectant toutes les structures ?}\footnote{C'est la
question des anneaux catégoriels, ou 2-anneaux. Sous la traduction employée
tout au long de cette lecture, elle demande la classification des spectres en
anneaux $E_{\infty}$ tronqués en degré $1$ — le fait que $\pi_{0}$ soit un
anneau et $\pi_{1}$ un module en étant la part immédiate, et le reste, comme
la page le pressent, ne l'étant pas.} C'est le dernier mot mathématique du
dossier.

\end{document}
